\documentclass[10pt,a4paper]{amsart}
\usepackage[margin=19mm]{geometry}
\usepackage{amsmath,amssymb,mathtools}
\usepackage[scr=rsfs,cal=euler]{mathalfa}
\usepackage{xfrac}
\usepackage[unicode,hidelinks,bookmarks=false]{hyperref}
\newcommand{\ie}{i.e.\ }
\newcommand{\cf}{cf.\ }
\newcommand{\resp}{respectively\ }
\newcommand{\Subsection}{\subsection}
\providecommand{\nobreakdash}{\nobreak\hbox{-}}
\setlength{\emergencystretch}{2em}
\multlinegap0pt
\usepackage{bm}
\usepackage{bbm}
\usepackage{dsfont}
\usepackage{xspace}
\usepackage{cancel}
\usepackage{mathtools}
\usepackage{amsthm}
\makeatletter
\@ifundefined{proposition}{\newtheorem{proposition}{Proposition}}{}
\makeatother
\title[On the Complexity of Counting Orderings in Graphs]{On the Complexity of Counting Orderings in Graphs}
\author{Marcelo Arenas, Mar\'ia Alejandra Schild, Bernardo Subercaseaux}
\date{Source: arXiv:2606.29157v1 (CC BY 4.0)}
\begin{document}
\maketitle

\noindent\emph{Source and licence.} On the Complexity of Counting Orderings in Graphs. Version arXiv:2606.29157v1 (CC BY 4.0), distributed under the Creative Commons Attribution 4.0 International licence, as recorded in the arXiv licence metadata for this exact version and shown on the abstract page https://arxiv.org/abs/2606.29157v1. Full rights notice: licenses/arxiv-2606.29157-CC-BY-4.0.txt.\par\smallskip

\noindent\textit{Editorial scope.} Counted unit: the Proposition whose source label is \texttt{prop:compute-polynomials} in the article (source0.tex lines 435--439, source label \texttt{prop:compute-polynomials}), delivered together with the article's own complete original proof of it (source0.tex lines 441--502, one \texttt{proof} environment of 62 lines). No mathematics is written, repaired or re-derived: statement and proof are the article's own text line for line. Required context retained uncounted: none. Every definition, display and label of those lines is retained unchanged. Omitted from this package and not redistributed: 1--434, 440--440, 503--1442. Nothing inside a retained excerpt is omitted or altered: every retained body is delivered exactly as the article's lines are written, so no omission receipt of any kind accompanies this package. Citations: the retained excerpts carry no citation command at all, so no bibliography entry of the article is transcribed and no reference is redistributed. Reference adaptation: none was needed and none was made; every reference inside the delivered unit resolves inside it, so no unresolved reference or citation is produced. Printed slips kept literal: no printed word, symbol, hypothesis or number of the counted statement or of its proof was altered. Layout and class: the article's own class is \texttt{amsart} with packages including preamble; the wrapper is the lane's pinned \texttt{amsart} editorial wrapper at 10pt on A4, which restates the theorem-like environments the retained text needs and adds the article's own macro definitions that the retained text needs (none), with extra wrapper packages (bm, bbm, dsfont, xspace, cancel, mathtools). The delivered numbering is the unit's own. Assets: no retained span contains a figure environment, a graphics-inclusion command, a PDF-page inclusion, a table or a TikZ picture; no image file is copied into this package, the package declares no asset dependency and no image file is redistributed with it. Licence: version arXiv:2606.29157v1 of this article is distributed under the Creative Commons Attribution 4.0 International licence, as recorded in the arXiv licence metadata for this exact version and shown on the abstract page https://arxiv.org/abs/2606.29157v1. A full rights notice is delivered at licenses/arxiv-2606.29157-CC-BY-4.0.txt. Characters: every retained excerpt is pure ASCII, so the delivered engine, XeLaTeX with its default Latin Modern text font, reports no missing character.
\begin{proposition}\label{prop:compute-polynomials}
   Given an instance $I = (\mathcal{S}, U)$ of $\#\mathsf{X3C}$, and an oracle for $\#\mathsf{SVO}(\cdot;\cdot)$, we can compute in polynomial time (in terms of the size of the instance $I$) two polynomials $R$ and $D$ in canonical form such that the following equality between rational functions holds:
   \[ \sum_{\sigma \in [n]!} \chi_{\sigma}(q) = \frac{R(q)}{D(q)}.
    \]
\end{proposition}

\begin{proof}
%     First, to each vertex $v \in V(G_{I, q}) \setminus \{r\}$, we associate uniformly at random a ``time'' $t(v) \in [0, 1]$. 
%     Letting $N := |V(G_{I, q})| - 1$, we will treat $t$ as a random variable in $[0,1]^N$. 
%     Note that $t$ naturally induces an ordering of the vertices after $r$: $\pi_t = (v_1, \dots, v_N)$ such that $t(v_1) < t(v_2) < \dots < t(v_N)$.\footnote{Since ties have probability $0$, we can safely ignore them.} 
%     Because this process is completely symmetric among all vertices, each vertex ordering occurs with probability $(N!)^{-1}$.  

%     Now, for an element $u$ in the universe $U(I)$, let $M_t(u)$ be the random variable corresponding to the smallest $t(v)$ s.t. $\{u,v\}$ is an edge in $G_{I,q}$. Equivalently, since by construction such $v$ must be one of the $S_i$ vertices coming from $\mathcal{S}(I)$, we can think of $M_t(u)$ as the earliest time assigned to a set containing $u$.
%      Then, observe that for an ordering of the vertices to be an SVO, we need that for each element $u \in U(I)$, all its $q$ copies $u_i$ hold $t(u_i) > M_t(u)$. Thus, letting $t(\mathcal{S}) \in [0,1]^{|\mathcal{S}|}$ denoting the random variable corresponding to the times only for the vertices in $\mathcal{S}$ and similarly for $t(U)$, we have 
%      \begin{align*}
%         P_I(q) &= \Pr_{t(U), t(\mathcal{S})}\left[\pi_t \text{ is an SVO} \right] \\ 
%         &= \Pr_{t(U), t(\mathcal{S})}\left[\bigcap_{u \in U(I), i \in [q]} \{t(u_i) > M_t(u)\}\right] \\ &= \mathbb{E}_{t(\mathcal{S})}\left[\Pr_{t(U)}\left[ \bigcap_{u \in U(I), i \in [q]} \{t(u_i) > M_t(u)\} \,\Bigg\vert\, t(\mathcal{S}) \right] \right].
%      \end{align*}
%      Since conditioned on $t(\mathcal{S})$, the value of $M_t(u)$ is fixed for each $u \in U(I)$, we obtain
%      \[
%          P_I(q) =  \int_{[0,1]^{|\mathcal{S}|}} \prod_{u \in U(I)} (1 - M_t(u))^q \; \mathrm{d}t.
%      \]
%      Now, up to a set of measure $0$, we can write $[0,1]^{|\mathcal{S}|} = [0,1]^n$ as the disjoint union of $n!$ simplices of the form \[ \Delta_\sigma := \{ 0 < t(S_{\sigma(1)}) < t(S_{\sigma(2)}) < \dots < t(S_{\sigma(n)}) < 1\},\] one for each permutation $\sigma$ of $[n]$.
%      Under this perspective, for a given permutation $\sigma$ and index $j \in [n]$ let 
%      \[
%         c_j(\sigma) := \left| S_{\sigma(j)} \setminus \bigcup_{i=1}^{j-1} S_{\sigma(i)} \right|. 
%      \]
%      In other words, $c_j(\sigma)$ is the number of elements of $U(I)$ whose earliest covering set is $S_{\sigma(j)}$.
%      Then, note that $0 \leq c_j(\sigma) \leq 3$, and $\sum_{j=1}^n c_j(\sigma) = |U(I)|$.
%      Furthermore, for a given $\sigma$, over the simplex $\Delta_\sigma$ we have 

%      \[
%          \prod_{u \in U(I)} (1 - M_t(u))^q  = \prod_{i=1}^n (1 - t(S_{\sigma(i)}))^{q \cdot c_i(\sigma)}.
%      \]
%      Therefore, we can rewrite $P_I(q)$ as
%      \[
%         P_I(q) = \sum_{\sigma \in [n]!} \int_{\Delta_\sigma} \prod_{i=1}^n (1 - t(S_{\sigma(i)}))^{q \cdot c_i(\sigma)} \; \mathrm{d}t.
%      \]

%      Now, we will need the following fact regarding such simplex integrals. 
%      \begin{lemma}\label{lemma:integrals}
%     Let $n \in \mathbb{N}_{\geq 1}$, and let $\beta_1, \dots, \beta_n$ be arbitrary non-negative reals. Then,
%     \[
%         \int_{0 < y_1 < \dots < y_n < 1} \prod_{i=1}^n y_i^{\beta_i} \mathrm{d}\vec{y} = \prod_{i=1}^n \frac{1}{ i +\sum_{j=1}^i \beta_j}.
%     \]
% \end{lemma}
% \noindent

% By \Cref{lemma:svo-formula} we know that

% \[
%         \frac{\#\mathsf{SVO}(G_{I, q}; r)}{(n+3hq)!} = \sum_{\sigma \in [n]!} \chi_{\sigma}(q),
% \]
% where
% \[
%     \chi_{\sigma}(q) \coloneqq  \frac{1}{\prod_{i=1}^n \left(i + \sum_{j=1}^i q \cdot c_{n-j+1}(\sigma)\right)}.
% \]

Note that all terms $i + \sum_{j=1}^i q \cdot c_{n-j+1}(\sigma)$ in the denominator of $\chi_{\sigma}(q)$ are of the form $Cq + \ell$ for integers $0 \leq C \leq |U|$ and $1 \leq \ell \leq n$. Thus, 
the polynomial 
\[
    D(q) := \prod_{C = 0}^{|U|} \prod_{\ell=1}^{n} (Cq + \ell)^n
\]
is necessarily divisible by the polynomial $\frac{1}{\chi_\sigma(q)}$ for every $\sigma$, making $D(q) \cdot \chi_\sigma(q)$ a polynomial of degree at most $\deg D \leq |U|n^2$. Therefore,
\[
    \sum_{\sigma \in [n]!} D(q)\cdot \chi_\sigma(q) = D(q) \cdot \sum_{\sigma \in [n]!} \chi_\sigma(q)
\] is also a polynomial, $R(q)$, of degree at most $|U| \cdot n^2$.  Thus, we can obtain its coefficients by first evaluating $\frac{\#\mathsf{SVO}(G_{I, q}; r)}{(n+3hq)!}$ at the points $1, 2, \ldots, |U|n^2 + 1$ (via the oracle) and $D$ at the same points (directly, since we know $D$ explicitly), then applying Lagrange interpolation, and finally putting the result into canonical form. Notice that $D$ can also be put into canonical form in polynomial time.
\end{proof}

\end{document}
